\begin{textsample}{POS Dim 1 – ro – Score 54.00 – ro/ro\_em\_11.txt}  \label{ex:f1_pos_143}
1.5.2.
O Jovem e o seu Projeto de Vida A \textbf{humanidade} e a sociedade contribuem para \textbf{que} as pessoas sejam seres de relações se tornem sociáveis.
O ser humano é \textbf{um} ser aberto ao diálogo e ao convívio social.
Essa forma de vivência em sociedade tem seus primeiros ensinamentos no âmbito familiar, onde os pais devem orientar seus filhos sobre o \textbf{que} é \textbf{justo} e correto, sobre o respeito ao próximo, entre outros.
Após a família, surge a Igreja como fonte de ensinamentos religiosos e espirituais.
Em seguida, destaca -se o papel da escola, \textbf{que} além de transmitir conhecimentos aos estudantes, deve orientá -los sobre o projeto de conhecer a si mesmo, de acolher o outro, sobre ensinamentos éticos, sobre o papel de cada \textbf{um} frente à sociedade, sobre as escolhas pessoais e profissionais estabelecidas longo de sua vida.
Em relação às escolhas, pelo fato de \textbf{que} os seres humanos \textbf{vivem} em sociedade, podem se tornar coletiva, \textbf{uma} vez \textbf{que} escolher para si é, também, escolher para o outro.
Como exemplo, pode -se destacar a situação em \textbf{que} \textbf{um} jovem toma \textbf{uma} decisão e conquista algo \textbf{que} agrada sua família e as pessoas \textbf{que} fazem parte do seu convívio social, \textbf{logo} a conquista se \textbf{revela} como sendo de caráter coletivo. \textbf{Contudo}, caso a tomada de decisão e a escolha não sejam realizadas de forma intencional, autônoma e crítica, o estudante poderá se perceber diante de \textbf{um} cenário em \textbf{que} o mesmo não tem consciência de escolher aquilo \textbf{que} é melhor para si, como para o outro.
Nessa perspectiva, fazer escolhas é projetar a vida, é olhar para frente, para além do \textbf{horizonte}, para além daquilo \textbf{que} se pode \textbf{ver}, e sonhar aquilo \textbf{que} se quer alcançar.
E isso não é \textbf{uma} \textbf{tarefa} fácil : \textbf{necessita} de \textbf{compromisso} sério consigo mesmo e com o outro.
Os jovens possuem \textbf{inúmeros} desejos, sonhos, metas \textbf{que} buscam alcançar.
No entanto, não é possível permanecer somente nos sonhos : deve -se lutar e buscar com afinco aquilo \textbf{que} se quer.
Conquistar seus objetivos é o desejo de qualquer pessoa, independentemente da idade \textbf{que} se possa ter.
Sendo assim, é relevante considerar \textbf{que} a grande maioria dos jovens tem aspirações ativas e/ou adormecidas \textbf{que} perpassam pela necessidade de superação, de ir além e conquistar tudo aquilo \textbf{que} desejam para sua vida pessoal, social e/ou profissional.
Diante de tal cenário, muitas vezes, o jovem não sabe qual o caminho \textbf{que} deve seguir.
E saber o itinerário a percorrer, é fundamental dentro da proposta do Novo Ensino Médio, pois o estudante é considerado como protagonista de sua própria formação e de sua história de vida.
Para tanto, no \textbf{escopo} dessa nova proposta de ensino, o Projeto de Vida oportunizará \textbf{uma} \textbf{direção}, \textbf{um} norte, para \textbf{que} os estudantes possam se orientar perante às \textbf{inúmeras} possibilidades e espaços existentes na sociedade para \textbf{que} o mesmo faça escolhas e projete sua vida. \textbf{Um} exemplo disso é \textbf{que}, após a conclusão do Ensino Médio, o estudante pode optar entre ir para o mercado de trabalho ou continuar seus estudos no Ensino Superior, ou até mesmo fazer ambas as coisas : estudar e trabalhar. \textbf{Contudo}, essa tomada de decisão e escolha, também requer indagações internas ao estudante referente à dimensão da sua vida pessoal e os caminhos \textbf{que} pretende trilhar para vida em sociedade com \textbf{compromisso} ético.
Partindo dessa premissa, a decisão do jovem deve perpassar necessariamente pela projeção no presente de \textbf{um} futuro pautado nas multidimensões pessoal, social, cidadã e profissional.
Com base nisso, o Projeto de Vida poderá auxiliar os estudantes nesses momentos de decisões em relação ao seu presente e ao seu futuro.
Até porque, “ quando se fala de Projeto de Vida, o trabalho intencional e reflexivo junto ao estudante em relação aos objetivos, \textbf{ideais} e sonhos \textbf{que} o acompanham, \textbf{exige} a organização, o planejamento e a estruturação de metas em curto, médio e longo prazo \textbf{que} podem envolver tanto o presente, como o futuro. ” Dessa forma, a fase da escolha não é \textbf{uma} \textbf{tarefa} fácil, pois requer tomadas de decisões sobre situações presentes nas dimensões da vida pessoal, social, cidadã e profissional.
Por esse motivo, ter \textbf{um} projeto de vida estruturado, possibilitando aos jovens tomada decisões sobre aquilo \textbf{que} querem conquistar e ser no futuro e/ou até mesmo escolher o \textbf{que} projetará para o seu futuro com o intuito buscar a concretude daquele projeto, \textbf{revela} -se como sendo desafiador ao estudante e ao professor \textbf{que} atuará como mediador nesse processo de construção.
Sob a perspectiva da concretude, é importante reconhecer \textbf{que} ao projetar, sonhar e desejar o futuro, o jovem se \textbf{vê} diante de \textbf{inúmeras} possibilidades e caminhos \textbf{que} poderão ser percorridos para alcançar suas metas.
Por essa razão, o ato de projetar a vida do ponto de vista abstrato contribui para vislumbrar o projeto de vida como algo incerto.
Diante disso, tomar decisões não é nada fácil, ainda mais quando elas, em sua grande maioria, têm a interferência e \textbf{influência} de outros, sendo estes, os pais ou os responsáveis legais, família, amigos.
Em relação à família, muitas vezes, os jovens se sentem coagidos a não decidirem por si mesmos, mas suas escolhas nada mais são do \textbf{que} acatar a \textbf{vontade} de outro(s ).
Assim, muitos jovens são \textbf{influenciados}, e até mesmo se sentem obrigados, a fazerem a \textbf{vontade} dos pais ou responsáveis como forma de não os desapontar.
Em \textbf{face} ao \textbf{exposto}, a partir do protagonismo juvenil na tomada de decisão e nas escolhas, o jovem avança o degrau da construção de si mesmo e do seu “ eu interior ”, entretanto essa construção não é finalizada ao final da Educação Básica, mas sim continuada no decorrer da sua vida.
Nesse sentido, \textbf{uma} das primeiras atividades desenvolvidas no projeto de vida é propiciar ao jovem o conhecimento de si mesmo.
Sobre o ato de conhecer seu “ Eu interior ”, cabe destacar não se trata de algo novo, pois entre os gregos antigos, o filósofo Sócrates ( \textbf{século} V a.C. ) já \textbf{revelava} o papel importante e os sentidos de as pessoas conhecerem a si mesmo. \textbf{Uma} das frases \textbf{conhecidas} \textbf{que} esse filósofo aborda consiste em “ conhece -te a ti mesmo ”.
E, assim, partindo do \textbf{método} chamado “ ironia ” e “ maiêutica ”, no campo filosófico, é possível conduzir as pessoas a olharem para dentro de si e reconhecer o quão desprovido de conhecimento pessoal e intelectual podem estar.
Diante dessas premissas, será \textbf{que} os jovens já estão preparados para fazer essas escolhas?
Será \textbf{que} eles sabem, de fato, aquilo \textbf{que} querem ser e aquilo \textbf{que} querem conquistar?
Ao refletir sobre tais questionamentos, vislumbra -se \textbf{que} a escolha exerce \textbf{um} papel de \textbf{suma} importância na sociedade, já \textbf{que} propiciar ao estudante a compreensão do princípio de partida pelo presente e o planejamento do futuro, requer o entendimento e a distinção das \textbf{inúmeras} dimensões da vida em sociedade.
Sob esse ponto de vista, o jovem, de forma intencional, desenvolverá a consciência daquilo \textbf{que} quer ser, mostrando -se preparado para assumir seu papel social com responsabilidade e zelo.
De modo específico, no contexto do trabalho com o Projeto de Vida, quando o jovem se conhece e se projeta, ele vai se lançar na sociedade, vai assumir a sua função dentro do convívio social.
Para tanto, o ato educativo intencional com práticas \textbf{que} orientem o jovem no desempenho do seu papel perante a sociedade, oportunizando espaços de aprendizagem deve compreender necessariamente a \textbf{inter-relação} entre as dimensões da vida pessoal, social, profissional e cidadã como sendo \textbf{uma} premissa essencial para formação integral desse estudante.
Se escolher já é \textbf{uma} \textbf{tarefa} difícil para o jovem, conhecer a si mesmo torna -se algo \textbf{bastante} complexo, \textbf{uma} vez \textbf{que} o conhecimento de si mesmo é \textbf{uma} atividade fundamental para a dimensão pessoal.
Sendo assim, partindo do âmbito filosófico e psicológico \textbf{que} \textbf{permeia} o trabalho pedagógico com o Projeto de Vida, constitui \textbf{fundamento} básico a construção do “ Eu interior ”, ou seja, primeiramente conhecer quem você é, em segundo, planejar a si mesmo numa perspectiva de curto, médio e longo prazo.
A construção interior e exterior desse processo \textbf{exige} o desenvolvimento de \textbf{um} trabalho pedagógico contínuo e integrado com todos os agentes da comunidade escolar.
Além disso, é relevante mencionar \textbf{que} tal processo não se esgota no ambiente escolar, bem como não terá fim, pois o \textbf{homem} nunca estará pronto, acabado.
Em suma, arquitetar o “ Eu ” e projetar -se diante das possibilidades \textbf{que} estão na frente de cada ser, \textbf{revela} -se como sendo complexo e desafiador. \textbf{Logo}, dentro da proposta educacional do Projeto de Vida de Educação Integral, o acolhimento é \textbf{um} fator \textbf{imprescindível} para \textbf{que} os novos estudantes se sintam acolhidos pela comunidade escolar e comecem a fazer parte dessa forma de ensino. \textbf{Uma} das premissas da escola deve ser conduzir esses recém-chegados no Ensino Médio à \textbf{percepção} da relevância da \textbf{aplicabilidade} do Projeto de Vida para o desenvolvimento integral e o planejamento de \textbf{um} futuro com vistas a traçar objetivos, metas e concretização de sonhos \textbf{que} desejam alcançar.
Com base nisso, o educador Antônio Carlos Gomes da Costa ( 2006 ) afirma : > Se ele [ \textbf{educando} ] tiver a sensação de \textbf{que} tem valor para alguém e de \textbf{que} é compreendido e aceito, vai olhar o futuro sem medo, será capaz de plasmar, de construir \textbf{um} projeto de vida.
Se ele constrói \textbf{um} projeto de vida, sua vida passa a ter \textbf{um} sentido ; se a vida passa a ter \textbf{um} sentido, ele começa a \textbf{ver} com outros olhos os estudos, a obediência, a profissionalização, o seguimento das regras, o tratamento com as pessoas etc.
Tudo isso se modifica na sua vida. ( COSTA, 2006, p. 69 ). \textbf{Face} ao \textbf{exposto}, o projeto de vida é \textbf{um} componente curricular \textbf{que} tem como objetivo auxiliar o estudante na construção de sua identidade pessoal.
Para tanto, é \textbf{crucial}, no desenvolvimento da identidade, conhecer a origem, etnia e a cultura do estudante. \textbf{Ademais}, é relevante ter conhecimento da base \textbf{conceitual} familiar do estudante, ao mesmo tempo em \textbf{que} se busca saber para onde o mesmo quer ir e aonde se quer chegar.
Para tanto, durante o processo de formação \textbf{identitária}, é necessário estabelecer critérios voltados para a dimensão pessoal, social, cidadã e profissional do ser fundamentado em princípios éticos e morais.
Cabe ressaltar \textbf{que} o projeto de vida é predominantemente \textbf{uma} construção de si mesmo, pois o indivíduo constrói sua própria identidade.
Ao se reconhecer como parte integrante de \textbf{um} \textbf{escopo} social, o jovem passa a exercer a sua cidadania e percebe \textbf{que} exerce \textbf{uma} função dentro do corpo social.
Em relação à abrangência dos tipos de Projeto de Vida, a \textbf{pesquisadora} professora Hanna Cebel Danza ( 2009 ) em sua pesquisa, identificou seis tipos de projetos de vida, sendo eles : conflito identitário, projeções normativas, projeções idealizadas, projeções parciais, projetos de vida e projetos de vida com \textbf{compromisso} ético.
E assim com vistas a ampliar a compreensão sobre os tipos de Projeto o organograma abaixo \textbf{ilustra} \textbf{uma} síntese sobre a visão da autora sobre cada projeto.
A partir da classificação dos tipos de Projetos de Vida ( Figura 2 ), percebe -se \textbf{que} é possível pensar no contexto escolar e nas estratégias para trabalhar o perfil de estudante a fim de mobilizar os saberes necessários para desenvolver o protagonismo autêntico, a tomada de decisão e autonomia do estudante frente ao ato de projetar sua vida.
Nesse sentido, tomando por base as características de cada tipo de projeto, percebe -se a complexidade do trabalho pedagógico com o Projeto de Vida no \textbf{que} \textbf{diz} respeito à diversidade presente na vida cotidiana dos estudantes e nos tipos de intervenções pedagógicas \textbf{que} poderão ser utilizadas para auxiliar na construção do seu projeto de vida individual de cada estudante.
Dessa forma, ao analisar o Projeto de conflito identitário ( figura 2 ), é pertinente destacar \textbf{que} o mesmo é \textbf{aqui} entendido como sendo o período de \textbf{crise} de identidade do estudante, \textbf{uma} vez \textbf{que} nesse tipo de projeto o jovem ainda não se encontrou e/ou ainda não sabe quem de fato é.
É importante destacar ainda \textbf{que} o perfil do jovem nesse tipo de projeto é reflexo direto de boa parte da juventude \textbf{que} não parou para pensar no amanhã e não desenvolveram ainda a pretensão em projetar e/ou idealizar o futuro.
Dentre as características principais, esse grupo de jovens geralmente concentra -se em \textbf{viver} o \textbf{hoje}, portanto, seu grande desafio é conhecer a si mesmo e criar sua própria identidade.
Frente a tal conflito identitário, cabe \textbf{salientar} \textbf{que} não envolve somente o fator pessoal, mas é também o fator social. \textbf{Logo}, cabe ao professor o desafio de mobilizar saberes e buscar estratégias para \textbf{que} esse jovem possa se projetar para o futuro e encontrar o seu lugar na sociedade a fim de buscar a concretude do Projeto de Vida.
Com relação ao “ Projeto com Projeções normativas ” ( figura 2 ), inicialmente é relevante esclarecer \textbf{que} o termo “ projeções ” \textbf{remete} ao entendimento de \textbf{que} ainda não existe \textbf{um} projeto de vida concretizado, \textbf{logo}, percebe -se \textbf{que} o jovem já carrega consigo intenções \textbf{baseadas} naquilo \textbf{que} acredita e nas suas expectativas sociais.
Do ponto de vista do aprofundamento, a intencionalidade do jovem em projetar \textbf{um} projeto com sob a perspectiva da projeção normativa enfraquece as multidimensões do ato de projetar sua vida, \textbf{uma} vez \textbf{que} nesse tipo de projeto, a escolha do jovem pautou -se simplesmente em \textbf{viver} a vida da forma como ele acredita.
Como exemplo disso, o jovem pode na tomada de decisão, escolher se quer trabalhar, casar e ter \textbf{uma} família, e ficar somente nisso o seu projeto de vida.
Entretanto, para essa situação específica, a intencionalidade da escolha do estudante, possibilita \textbf{um} projeto de vida vago sem \textbf{fundamento}, sem valores pessoais e sociais.
Em se tratando do tipo de projeto de vida, denominado “ projeções idealizadas ”, esse envolve o jovem \textbf{que} ainda não possui \textbf{um} projeto consolidado e \textbf{que} traz consigo ideias para o futuro, desejos, aspirações e sonhos de realização pessoal \textbf{que} ainda se encontram desvinculados de situações concretas.
Sendo assim, o jovem ao ingressar no Ensino Médio apresenta \textbf{uma} projeção idealizada de vida e constrói ideias, mas não desenvolve \textbf{um} planejamento pessoal para \textbf{que} possa consolidar suas projeções.
Em \textbf{face} disso, esse jovem deseja \textbf{inúmeras} coisas para o seu futuro, no entanto, as projeções permanecem na esfera da intenção abstrata e totalmente desvinculadas à realidade \textbf{vivida} desse estudante e das ações \textbf{que} o mesmo realiza para conquistar tais objetivos e metas.
No \textbf{que} \textbf{diz} respeito ao tipo de “ Projeto Parcial ”, constatam -se como características \textbf{centrais} \textbf{que} o jovem é capaz de projetar sua vida, de forma \textbf{bastante} aprofundada, mas essa projeção é parcial.
Para tanto, nesse tipo de projeto percebe -se a mobilização de saberes do estudante \textbf{que} projeta com \textbf{bastante} especificidade a importância da vida familiar, entretanto, encontra limitações para projetar com a mesma clareza sua vida profissional e cidadã.
O jovem \textbf{que} \textbf{consegue} construir \textbf{um} Projeto Parcial, ele \textbf{consegue} \textbf{argumentar} com profundidade sobre sua família, bem como evidencia o sentido, a importância e os valores morais \textbf{que} a mesma exerce na sua vida e na sociedade. \textbf{Contudo}, na esfera social, profissional e cidadã não se sente preparado para deliberar sobre a importância do mundo do trabalho em sua vida e no campo social.
Em relação à classificação “ Projeto de vida ”, é relevante mencionar \textbf{que} contempla significação pessoal e está vinculado a ações concretas do ato de projetar a vida no presente e no futuro.
O jovem capaz de construir esse tipo de Projeto desenvolve competências para projetar sua vida numa perspectiva de curto, médio e longo prazo.
Por fim, “ o Projeto de Vida com \textbf{compromisso} ético ”, caracteriza -se como sendo o foco \textbf{central} do trabalho pedagógico com o jovem da \textbf{contemporaneidade}, tendo em vista a superação das demandas sociais e os conflitos éticos presente no cenário social.
Nesse sentido, \textbf{instigar} o jovem a projetar sua vida com responsabilidade social, ambiental e ética, deve ser considerada a primordial fonte de inspiração para o projeto de vida desejado ao estudante do Ensino Médio do Estado de Rondônia.

% matched lemmas: ademais, aplicabilidade, aqui, argumentar, basear, bastante, central, compromisso, conceitual, conhecido, conseguir, contemporaneidade, contudo, crise, crucial, direção, dizer, educar, escopo, exigir, exposto, face, fundamento, hoje, homem, horizonte, humanidade, ideal, identitária, ilustrar, imprescindível, influenciar, influência, instigar, inter-relação, inúmero, justo, logo, método, necessitar, percepção, permear, pesquisador, que, remeter, revelar, salientar, sumo, século, tarefa, um, ver, viver, vontade
\end{textsample}
